\documentclass[11pt,a4paper]{article}
\usepackage[utf8]{inputenc}
\usepackage[czech]{babel}
\usepackage[T1]{fontenc}
\usepackage{geometry}
\usepackage{enumitem}
\usepackage{xcolor}
\usepackage{tcolorbox}
\usepackage{booktabs}
\usepackage{array}
\usepackage{hyperref}
\usepackage{fancyhdr}
\usepackage{titlesec}
\usepackage{amssymb}

\geometry{margin=2cm}
\pagestyle{fancy}
\fancyhf{}
\rhead{Objektové programování 2 -- UAI/520}
\lhead{Příprava na test}
\cfoot{\thepage}

\definecolor{darkblue}{RGB}{0,51,102}
\definecolor{lightblue}{RGB}{230,240,250}
\definecolor{warningred}{RGB}{180,30,30}
\definecolor{successgreen}{RGB}{30,120,30}
\definecolor{highlightyellow}{RGB}{255,250,205}

\tcbset{
    defbox/.style={colback=lightblue,colframe=darkblue,fonttitle=\bfseries,boxrule=1pt},
    warnbox/.style={colback=red!5,colframe=warningred,fonttitle=\bfseries,boxrule=1pt},
    tipbox/.style={colback=highlightyellow,colframe=orange!80!black,fonttitle=\bfseries,boxrule=1pt}
}

\titleformat{\section}{\Large\bfseries\color{darkblue}}{\thesection}{1em}{}
\titleformat{\subsection}{\large\bfseries\color{darkblue!80}}{\thesubsection}{1em}{}
\titleformat{\subsubsection}{\normalsize\bfseries\color{darkblue!60}}{\thesubsubsection}{1em}{}

\begin{document}

\begin{center}
{\Huge\bfseries\color{darkblue} Příprava na test}\\[0.5em]
{\Large Objektové programování 2 -- UAI/520}\\[0.3em]
{\normalsize Zpracováno z prezentací a vzorového testu Verity}
\end{center}

\tableofcontents
\newpage

%=============================================================================
\section{Práce s chybami a výjimky}
%=============================================================================

\subsection{Základní terminologie}

\begin{tcolorbox}[defbox,title=Defenzivní programování]
Předvídání, že se věci mohou pokazit. Server by měl předpokládat, že klienti jsou \textbf{potenciálně nepřátelští}, ne že se chovají dobře.
\end{tcolorbox}

\textbf{Typické chybové situace:}
\begin{itemize}[leftmargin=*]
    \item Nesprávná implementace (nevyhovuje specifikaci)
    \item Nepatřičná žádost o objekt (neplatný index)
    \item Nekonzistentní nebo nevhodný stav objektu
    \item Problémy dané prostředím: nesprávné URL, přerušení sítě, chybějící soubory, nedostatečná oprávnění
\end{itemize}

\subsection{Kategorie výjimek}

\begin{center}
\begin{tabular}{>{\raggedright}p{3.5cm} >{\raggedright}p{5cm} >{\raggedright\arraybackslash}p{5cm}}
\toprule
\textbf{Typ} & \textbf{Použití} & \textbf{Příklad} \\
\midrule
\textbf{Kontrolované} (checked) & Podtřída \texttt{Exception}. Předvídatelná selhání, zotavení možné. & \texttt{IOException}, \texttt{FileNotFoundException} \\
\addlinespace
\textbf{Nekontrolované} (unchecked) & Podtřída \texttt{RuntimeException}. Nepředvídatelná selhání, zotavení nepravděpodobné. & \texttt{IllegalArgumentException}, \texttt{NullPointerException} \\
\bottomrule
\end{tabular}
\end{center}

\begin{tcolorbox}[warnbox,title=Klíčové rozdíly]
\begin{itemize}[leftmargin=*]
    \item \textbf{Kontrolované}: Kompilátor \textbf{vynucuje} zpracování (try-catch nebo throws). Používají se pro \textbf{externí} problémy (soubory, síť).
    \item \textbf{Nekontrolované}: Kompilátor \textbf{nekontroluje}. Používají se pro \textbf{programátorské} chyby. Způsobí ukončení programu, pokud nejsou odchyceny.
\end{itemize}
\end{tcolorbox}

\subsection{Vyhazování výjimek}

\begin{verbatim}
public ContactDetails getDetails(String key) {
    if(key == null) {
        throw new IllegalArgumentException("null key in getDetails");
    }
    if(key.trim().length() == 0) {
        throw new IllegalArgumentException("Empty key passed to getDetails");
    }
    return book.get(key);
}
\end{verbatim}

\textbf{Důležité:}
\begin{itemize}
    \item Výjimka přeruší normální tok řízení -- metoda se předčasně ukončí, nevrací hodnotu
    \item Klient může výjimku odchytit (catch) nebo nechat propagovat
    \item Dokumentace: používej \texttt{@throws} v Javadoc
\end{itemize}

\subsection{Zpracování výjimek -- try-catch-finally}

\begin{verbatim}
try {
    addressbook.saveToFile(filename);
    successful = true;
} catch(IOException e) {
    System.out.println("Unable to save to " + filename);
    successful = false;
} finally {
    // Provede se VŽDY -- i při return v try/catch
    // Uzavření souborů, uvolnění zdrojů
}
\end{verbatim}

\textbf{Vícenásobné odchycení:}
\begin{verbatim}
try { ... }
catch(EOFException | FileNotFoundException e) { ... }  // Java 7+
\end{verbatim}

\subsection{Definice vlastních výjimek}

\begin{verbatim}
public class NoMatchingDetailsException extends Exception {
    private String key;
    public NoMatchingDetailsException(String key) { this.key = key; }
    public String getKey() { return key; }
    public String toString() {
        return "No details matching '" + key + "' were found.";
    }
}
\end{verbatim}

\begin{tcolorbox}[tipbox,title=Pravidlo pro výběr nadtypu]
\begin{itemize}
    \item Rozšiř \texttt{RuntimeException} $\rightarrow$ nekontrolovaná
    \item Rozšiř \texttt{Exception} $\rightarrow$ kontrolovaná
\end{itemize}
\end{tcolorbox}

\subsection{Aserce (Assertions)}

\begin{tcolorbox}[warnbox,title=Assertions vs. Výjimky]
\textbf{NEPLATÍ}: Aserce nejsou alternativou k výjimkám!
\begin{itemize}
    \item Aserce: \textbf{interní kontroly} konzistence, používají se ve vývoji, odstraňují se z produkce
    \item Výjimky: \textbf{externí} problémy, zůstávají v produkčním kódu
\end{itemize}
\end{tcolorbox}

\begin{verbatim}
assert !keyInUse(key);  // kontrola po removeDetails
assert consistentSize() : "Nekonzistentní velikost adresáře";
\end{verbatim}

\textbf{Pozor}: \texttt{assert} se vypíná v produkci (\texttt{-ea} flag). \textbf{Nezahrnuj} do assert normální funkcionalitu!

\subsection{Zotavení z nestandardních stavů}

\begin{verbatim}
boolean successful = false;
int attempts = 0;
do {
    try {
        contacts.saveToFile(filename);
        successful = true;
    } catch(IOException e) {
        attempts++;
        if(attempts < MAX_ATTEMPTS) {
            filename = alternativniJmenoSouboru;
        }
    }
} while(!successful && attempts < MAX_ATTEMPTS);
\end{verbatim}

\textbf{Strategie vyhýbání se selháním:}
\begin{itemize}
    \item Používej dotazovací metody serveru (např. \texttt{keyInUse()}) před voláním rizikové operace
    \item Nekontrolované výjimky pro skutečně neočekávané stavy
\end{itemize}

%=============================================================================
\section{Genericita (Generics)}
%=============================================================================

\subsection{Základní koncepty}

\begin{tcolorbox}[defbox,title=Terminologie genericity]
\begin{tabular}{ll}
\textbf{Typový parametr} & Zástupný znak (např. \texttt{T}, \texttt{E}, \texttt{K}, \texttt{V}) \\
\textbf{Typový argument} & Konkrétní typ dosazený za parametr \\
\textbf{Instance generického typu} & Generický typ po dosazení argumentu \\
\textbf{Raw typ} & Použití bez typového argumentu (\texttt{List} místo \texttt{List<String>}) \\
\end{tabular}
\end{tcolorbox}

\textbf{Jmenné konvence:}
\begin{center}
\begin{tabular}{cl}
\toprule
\textbf{Značka} & \textbf{Význam} \\
\midrule
E & Element (kolekce) \\
K & Key (klíč) \\
V & Value (hodnota) \\
T & Type (obecný typ) \\
N & Number \\
S, U, V & Další typové parametry \\
\bottomrule
\end{tabular}
\end{center}

\subsection{Přínosy genericity}

\begin{enumerate}
    \item \textbf{Vyšší vypovídací schopnost} -- kód je čitelnější
    \item \textbf{Eliminace přetypování} -- žádné \texttt{(Integer) list.get(0)}
    \item \textbf{Silnější typová kontrola při překladu} -- chyby se odhalí dříve
    \item \textbf{Možnost užití generických algoritmů} -- odstranění opakujícího se kódu
\end{enumerate}

\begin{tcolorbox}[warnbox,title=Nebezpečí raw typů]
\texttt{Box rawBox = new Box();} -- kompilátor neví o typu, varování při volání generických metod. \textbf{Může vést k chybám při běhu} (ClassCastException).
\end{tcolorbox}

\subsection{Meze typových parametrů}

\begin{verbatim}
// Horní omezení (extends)
public class NaturalNumber<T extends Integer> { ... }

// Vícenásobná omezení -- třída musí být první!
class D<T extends A & B & C> { ... }

// Wildcards
?                // libovolný typ (není Object!)
? extends Number // podtyp Number (čtení)
? super Integer  // nadtyp Integer (zápis)
\end{verbatim}

\begin{tcolorbox}[tipbox,title=PECS pravidlo]
\textbf{Producer Extends, Consumer Super}:
\begin{itemize}
    \item Čteš z kolekce? $\rightarrow$ \texttt{? extends T}
    \item Zapisuješ do kolekce? $\rightarrow$ \texttt{? super T}
\end{itemize}
\end{tcolorbox}

\subsection{Generické metody}

\begin{verbatim}
// Typový parametr před návratovým typem
public static <K, V> boolean compare(Pair<K, V> p1, Pair<K, V> p2) {
    return p1.getKey().equals(p2.getKey()) &&
           p1.getValue().equals(p2.getValue());
}

// Omezení na Comparable
public static <T extends Comparable<T>> int countGreaterThan(T[] arr, T elem) {
    int count = 0;
    for (T e : arr)
        if (e.compareTo(elem) > 0)
            ++count;
    return count;
}
\end{verbatim}

\subsection{Vymazání typu (Type Erasure)}

\begin{tcolorbox}[defbox,title=Co je type erasure?]
Překladač přepíše všechny typové parametry jejich mezemi nebo \texttt{Object} (pokud meze nejsou). V bajtkódu \textbf{nejsou} generické typy!
\end{tcolorbox}

\begin{verbatim}
// Zdrojový kód:
public class Node<T> {
    private T data;
    public T getData() { return data; }
}

// Po erasure:
public class Node {
    private Object data;  // nebo Comparable s extends
    public Object getData() { return data; }
}
\end{verbatim}

\textbf{Důsledky:}
\begin{itemize}
    \item Nelze vytvářet instance typových parametrů: \texttt{new E()} -- CHYBA
    \item Nelze deklarovat statické proměnné typového parametru
    \item Nelze vytvářet pole parametrických typů: \texttt{new List<Integer>[2]} -- CHYBA
    \item Nelze přetížit metody, které mají stejnou signaturu po erasure
\end{itemize}

\subsection{Wildcard capture}

\begin{verbatim}
// PROBLÉM: překladač neumí odvodit typ ?
void foo(List<?> i) {
    i.set(0, i.get(0));  // CHYBA: get vrací Object, set očekává ?
}

// ŘEŠENÍ: pomocná metoda
void foo(List<?> i) {
    fooHelper(i);
}
private <T> void fooHelper(List<T> l) {
    l.set(0, l.get(0));  // OK: T je odvozeno
}
\end{verbatim}

%=============================================================================
\section{Tvorba grafického rozhraní (Swing)}
%=============================================================================

\subsection{Architektura GUI}

\begin{tcolorbox}[defbox,title=Tři pilíře GUI]
\begin{enumerate}
    \item \textbf{Komponenty} -- stavební bloky (tlačítka, menu, posuvníky)
    \item \textbf{Rozvržení} -- manažeři rozvrhování (Layout Managers)
    \item \textbf{Události} -- reakce na uživatelský vstup
\end{enumerate}
\end{tcolorbox}

\textbf{Hierarchie balíčků:}
\begin{itemize}
    \item \texttt{java.awt} -- základní komponenty (AWT)
    \item \texttt{javax.swing} -- rozšířené komponenty (Swing) -- \textbf{používej tyto!}
\end{itemize}

\subsection{Struktura aplikačního okna}

\begin{center}
\begin{tabular}{cl}
\toprule
\textbf{Komponenta} & \textbf{Účel} \\
\midrule
\texttt{JFrame} & Hlavní aplikační okno \\
\texttt{JMenuBar} & Pruh menu pod titulkem \\
\texttt{JMenu} & Menu (např. File) \\
\texttt{JMenuItem} & Položka menu (např. Open) \\
\texttt{JPanel} & Kontejner pro seskupení komponent \\
\texttt{JLabel} & Textový/popiskový prvek \\
\texttt{JButton} & Tlačítko \\
\texttt{JTextField} & Vstupní textové pole \\
\bottomrule
\end{tabular}
\end{center}

\subsection{Manažeři rozvrhování (Layout Managers)}

\begin{center}
\begin{tabular}{>{\raggedright}p{2.5cm} >{\raggedright\arraybackslash}p{10cm}}
\toprule
\textbf{Layout} & \textbf{Chování} \\
\midrule
\texttt{FlowLayout} & Komponenty vedle sebe, zalamování jako text \\
\texttt{BorderLayout} & 5 oblastí: NORTH, SOUTH, EAST, WEST, CENTER \\
\texttt{GridLayout} & Mřížka s pevným počtem řádků a sloupců \\
\texttt{BoxLayout} & Vertikální nebo horizontální řada \\
\texttt{GridBagLayout} & Nejsložitější, pružné umístění \\
\bottomrule
\end{tabular}
\end{center}

\begin{tcolorbox}[tipbox,title=Vnořené kontejnery]
Pro sofistikované rozvržení použij \textbf{vnořené panely} (JPanel) s různými layouty. Často vhodnější než GridBagLayout!
\end{tcolorbox}

\subsection{Zpracování událostí}

\textbf{Přístupy:}

\begin{enumerate}
    \item \textbf{Centralizovaný} -- jeden objekt implementuje \texttt{ActionListener}, rozhoduje podle \texttt{getActionCommand()}
    \begin{itemize}
        \item Nevýhoda: špatně škáluje, závislost na textu komponenty
    \end{itemize}
    
    \item \textbf{Anonymní vnitřní třídy} -- každá komponenta má vlastní listener
    \begin{verbatim}
    openItem.addActionListener(new ActionListener() {
        public void actionPerformed(ActionEvent e) { openFile(); }
    });
    \end{verbatim}
    
    \item \textbf{Lambda výrazy} (Java 8+) -- nejčistší zápis
    \begin{verbatim}
    openItem.addActionListener(e -> openFile());
    \end{verbatim}
\end{enumerate}

\subsection{Vnořené třídy (Nested Classes)}

\begin{center}
\begin{tabular}{>{\raggedright}p{3cm} >{\raggedright}p{5cm} >{\raggedright\arraybackslash}p{5cm}}
\toprule
\textbf{Typ} & \textbf{Instance} & \textbf{Přístup k vnější třídě} \\
\midrule
Statická vnořená & Nezávislá na instanci vnější & Pouze statické členy \\
Vnitřní (inner) & Vázána na instanci vnější & Všechny členy, i private \\
Lokální vnitřní & V bloku metody & Všechny členy \\
Anonymní vnitřní & Bez hlavičky, jedna instance & Dědí z nadtypu/implements interface \\
\bottomrule
\end{tabular}
\end{center}

\begin{tcolorbox}[defbox,title=Anonymní vnitřní třída]
Deklarace bez jména, typ určen nadtypem. Používá se pro \textbf{jednorázovou} instanci listeneru.
\begin{verbatim}
// Syntaxe: new Supertyp() { ... implementace ... }
new ActionListener() {
    public void actionPerformed(ActionEvent e) { ... }
}
\end{verbatim}
\end{tcolorbox}

\subsection{Dialogy}

\begin{itemize}
    \item \textbf{Modální} -- blokují ostatní interakce (vyžadují odpověď)
    \item \textbf{Nemodální} -- umožňují další interakci (pozor na nekonzistence)
    \item \texttt{JOptionPane} -- standardní dialogy: Message, Confirm, Input
\end{itemize}

%=============================================================================
\section{JavaFX}
%=============================================================================

\subsection{Historie a dostupnost}

\begin{center}
\begin{tabular}{cl}
\toprule
\textbf{Rok} & \textbf{Verze} \\
\midrule
1996 & AWT (JDK 1.0) \\
1998 & Swing (J2SE 1.2) \\
2007 & JavaFX 1.0 (JavaFX Script) \\
2011 & JavaFX 2.0 (ukončení Script, čistá Java) \\
2014 & JavaFX 8 (součást Java 8) \\
2018+ & Java 9+: samostatné SDK (OpenJFX, Gluon) \\
\bottomrule
\end{tabular}
\end{center}

\subsection{Architektura JavaFX}

\begin{tcolorbox}[defbox,title=Hierarchie vrstev (odspoda nahoru)]
\begin{enumerate}
    \item \textbf{JVM} -- Java Virtual Machine
    \item \textbf{JDK API} -- základní knihovny
    \item \textbf{Prism} -- grafický engine (2D/3D, hardware acceleration)
    \item \textbf{Glass Windowing Toolkit} -- platformově závislá vrstva
    \item \textbf{Quantum Toolkit} -- spojuje Prism a Glass
    \item \textbf{JavaFX Public APIs + Scene Graph} -- vrstva pro vývojáře
\end{enumerate}
\end{tcolorbox}

\subsection{Klíčové rysy JavaFX}

\begin{itemize}
    \item FXML a Scene Builder -- deklarativní popis GUI
    \item CSS stylování -- vzhled oddělen od logiky
    \item WebView -- vkládání webových stránek (WebKit)
    \item Interoperabilita se Swingem (\texttt{SwingNode})
    \item 3D grafika, Canvas API, Printing API, Rich Text
    \item Multitouch, Hi-DPI, hardwarová akcelerace (Prism)
\end{itemize}

\subsection{Graf scény (Scene Graph)}

\begin{tcolorbox}[defbox,title=Struktura GUI JavaFX]
\begin{center}
\texttt{Stage} $\supset$ \texttt{Scene} $\supset$ \texttt{Root Pane} $\supset$ \texttt{Nodes}
\end{center}
\begin{itemize}
    \item \textbf{Stage} -- ``jeviště'', vnější obálka (okno aplikace)
    \item \textbf{Scene} -- ``scéna'', obsah okna
    \item \textbf{Root Pane/Container} -- \texttt{StackPane}, \texttt{GridPane}, \texttt{BorderPane}, \texttt{FlowPane}...
    \item \textbf{Nodes} -- tlačítka, texty, obrazce, 3D objekty
\end{itemize}
\end{tcolorbox}

\textbf{Hierarchie elementů (od nejvyšší):}
\begin{center}
\texttt{Stage} $>$ \texttt{Scene} $>$ \texttt{Pane/Container} $>$ \texttt{Controls/Nodes}
\end{center}

\begin{tcolorbox}[warnbox,title=Otázka ze vzorového testu!]
Který objekt je \textbf{nejvyše postavený} v hierarchii? $\rightarrow$ \textbf{Stage} (nikoli Scene, Pane ani GridLayout)
\end{tcolorbox}

\subsection{FXML a Controller}

\begin{tcolorbox}[defbox,title=Struktura FXML aplikace]
\begin{tabular}{ll}
\texttt{FXMLDocument.fxml} & Deklarativní popis GUI (generuje Scene Builder) \\
\texttt{FXMLDocumentController.java} & Logika interakce, anotace \texttt{@FXML} \\
\texttt{HelloWorldFXML.java} & Spouštěcí třída (\texttt{Application.launch()}) \\
\end{tabular}
\end{tcolorbox}

\begin{tcolorbox}[warnbox,title=Účel anotace @FXML]
\texttt{@FXML} slouží k \textbf{propojení identifikace komponent} mezi kontrolerem a definicí vzhledu GUI. Umožňuje injectovat komponenty z FXML do controlleru podle \texttt{fx:id}.

\textbf{Nejedná se o:} generování dokumentace, označení jazyka FXML.
\end{tcolorbox}

\begin{verbatim}
public class FXMLDocumentController implements Initializable {
    @FXML
    private Label label;  // injectnuto z FXML podle fx:id="label"
    
    @FXML
    private void handleButtonAction(ActionEvent event) {
        label.setText("Hello World!");
    }
    
    @Override
    public void initialize(URL url, ResourceBundle rb) {
        // Inicializace po načtení FXML
    }
}
\end{verbatim}

\subsection{Layouty v JavaFX}

\begin{center}
\begin{tabular}{>{\raggedright}p{3cm} >{\raggedright\arraybackslash}p{9cm}}
\toprule
\textbf{Layout} & \textbf{Použití} \\
\midrule
\texttt{BorderPane} & 5 oblastí: top, bottom, left, right, center \\
\texttt{HBox/VBox} & Horizontální/vertikální řada \\
\texttt{StackPane} & Vrstvení komponent na sobě \\
\texttt{GridPane} & Mřížka s definicí řádků/sloupců \\
\texttt{FlowPane} & Zalamování jako FlowLayout \\
\texttt{TilePane} & Uniformní dlaždice \\
\texttt{AnchorPane} & Kotvení k okrajům \\
\bottomrule
\end{tabular}
\end{center}

\subsection{Životní cyklus JavaFX aplikace}

\begin{enumerate}
    \item \texttt{launch(args)} -- statická metoda \texttt{Application}
    \item Vytvoření instance aplikace
    \item \texttt{init()} -- inicializace (mimo UI vlákno!)
    \item \texttt{start(Stage primaryStage)} -- hlavní metoda, vytvoření scény
    \item Běh aplikace -- zpracování událostí
    \item \texttt{stop()} -- ukončení (volá se při zavření)
\end{enumerate}

\begin{tcolorbox}[warnbox,title=Vlákna v JavaFX]
\begin{itemize}
    \item \texttt{init()} -- \textbf{spouštěcí vlákno} (ne UI!)
    \item \texttt{start()} a veškerá manipulace se scénou -- \textbf{JavaFX Application Thread (JAT)}
    \item Náročné operace -- \textbf{pracovní vlákna na pozadí} (\texttt{Task}, \texttt{Service})
\end{itemize}
\end{tcolorbox}

\subsection{Pracovní vlákna v JavaFX}

\begin{center}
\begin{tabular}{>{\raggedright}p{3cm} >{\raggedright\arraybackslash}p{9cm}}
\toprule
\textbf{Třída} & \textbf{Účel} \\
\midrule
\texttt{Task<T>} & Jednorázová úloha na pozadí, implementace \texttt{FutureTask} \\
\texttt{Service<T>} & Opakované spouštění Tasků, sledování stavu \\
\texttt{WorkerStateEvent} & Události při změně stavu (SUCCEEDED, FAILED...) \\
\bottomrule
\end{tabular}
\end{center}

\begin{verbatim}
Task<Integer> task = new Task<Integer>() {
    @Override
    protected Integer call() throws Exception {
        for (int i = 0; i < 100000; i++) {
            if (isCancelled()) break;
            updateProgress(i, 100000);  // bezpečná aktualizace UI
        }
        return i;
    }
};

ProgressBar bar = new ProgressBar();
bar.progressProperty().bind(task.progressProperty());
new Thread(task).start();
\end{verbatim}

%=============================================================================
\section{Vlákna v Javě (Concurrency)}
%=============================================================================

\subsection{Procesy vs. vlákna}

\begin{center}
\begin{tabular}{>{\raggedright}p{3cm} >{\raggedright}p{5cm} >{\raggedright\arraybackslash}p{5cm}}
\toprule
\textbf{Aspekt} & \textbf{Proces} & \textbf{Vlákno} \\
\midrule
Paměť & Vlastní paměťový prostor & Sdílí paměť procesu \\
Zdroje & Samostatná množina & Sdílí zdroje procesu \\
Komunikace & Sokety, potrubí & Přímý přístup k paměti \\
Náročnost & Vysoká & Nízká (lightweight) \\
\bottomrule
\end{tabular}
\end{center}

\subsection{Vytvoření a spuštění vlákna}

\textbf{Dva způsoby:}

\begin{enumerate}
    \item \textbf{Runnable} -- preferovaný, obecnější
    \begin{verbatim}
    public class HelloRunnable implements Runnable {
        public void run() { System.out.println("Hello from thread!"); }
        public static void main(String[] args) {
            new Thread(new HelloRunnable()).start();
        }
    }
    \end{verbatim}
    
    \item \textbf{Podtřída Thread} -- přímo dědit, překrýt \texttt{run()}
    \begin{verbatim}
    public class HelloThread extends Thread {
        public void run() { ... }
        // new HelloThread().start();
    }
    \end{verbatim}
\end{enumerate}

\begin{tcolorbox}[warnbox,title=Důležité!]
Vlákno se spouští metodou \texttt{start()}, nikoli \texttt{run()}! \texttt{run()} by jen provedlo kód v aktuálním vlákně.
\end{tcolorbox}

\subsection{Stavy vlákna}

\begin{center}
\begin{tabular}{cl}
\toprule
\textbf{Stav} & \textbf{Význam} \\
\midrule
\texttt{NEW} & Vytvořeno, ještě neběží \\
\texttt{RUNNABLE} & Připraveno ke běhu (může běžet, nemusí) \\
\texttt{BLOCKED} & Čeká na zámek (monitor) \\
\texttt{WAITING} & Čeká na notifikaci (\texttt{wait()}, \texttt{join()} bez timeoutu) \\
\texttt{TIMED\_WAITING} & Čeká s časovým limitem (\texttt{sleep()}, \texttt{join(ms)}) \\
\texttt{TERMINATED} & Ukončeno (normálně nebo výjimkou) \\
\bottomrule
\end{tabular}
\end{center}

\begin{tcolorbox}[warnbox,title=Otázka ze vzorového testu!]
Správná sekvence stavů: \textbf{NEW, Runnable, Blocked, Waiting, Timed waiting, Terminated}

\textbf{Pozor na}: Running není oficiální stav (je součástí Runnable), Stopped/Finished/Canceled nejsou standardní!
\end{tcolorbox}

\subsection{Přerušení vlákna (Interrupt)}

\begin{tcolorbox}[defbox,title=Mechanismus interrupt]
\textbf{Přerušení je žádost}, ne příkaz. Vlákno musí spolupracovat.
\begin{itemize}
    \item \texttt{Thread.interrupt()} -- nastaví příznak přerušení
    \item \texttt{Thread.interrupted()} -- zkontroluje a \textbf{vymaže} příznak
    \item \texttt{isInterrupted()} -- zkontroluje, \textbf{nevymaže}
\end{itemize}
\end{tcolorbox}

\begin{verbatim}
// Správná podpora přerušení:
public void run() {
    while (!Thread.interrupted()) {
        // práce...
        try {
            Thread.sleep(1000);  // vyhodí InterruptedException při interrupt
        } catch (InterruptedException e) {
            return;  // nebo break, nebo obnovení příznaku
        }
    }
}
\end{verbatim}

\subsection{Synchronizace}

\subsubsection{Problémy souběhu}

\begin{tcolorbox}[defbox,title=Thread Interference]
Když více vláken přistupuje ke sdíleným datům a provádí operace, které nejsou \textbf{atomické}. I jednoduchá operace jako \texttt{c++} se přeloží na: načti c, zvyš, ulož.
\end{tcolorbox}

\begin{tcolorbox}[defbox,title=Memory Consistency Errors]
Různá vlákna mají různý pohled na sdílenou paměť. Řešením je relace \textbf{happens-before} -- zaručená časová vazba mezi operacemi.
\end{tcolorbox}

\subsubsection{Nástroje synchronizace}

\begin{center}
\begin{tabular}{>{\raggedright}p{3.5cm} >{\raggedright\arraybackslash}p{9.5cm}}
\toprule
\textbf{Nástroj} & \textbf{Popis} \\
\midrule
\texttt{synchronized} metoda & Zámek na instanci (nebo třídě pro static). Jeden zámek na objekt -- více synchronizovaných metod se vzájemně blokují. \\
\addlinespace
\texttt{synchronized} blok & Explicitní specifikace objektu zámku. Umožňuje jemnější granularitu (různé zámky pro různá data). \\
\addlinespace
\texttt{ReentrantLock} & Explicitní zámek z \texttt{java.util.concurrent.locks}. Podporuje \texttt{tryLock()} (neblokující), podmínky (Condition). \\
\addlinespace
Atomické proměnné & \texttt{AtomicInteger}, \texttt{AtomicLong}, \texttt{AtomicReference} -- lock-free operace. \\
\bottomrule
\end{tabular}
\end{center}

\begin{verbatim}
// Synchronizované metody -- jeden zámek na instanci
public synchronized void increment() { c++; }
public synchronized void decrement() { c--; }

// Synchronizované bloky -- různé zámky pro nezávislá data
private Object lock1 = new Object();
private Object lock2 = new Object();
public void inc1() { synchronized(lock1) { c1++; } }
public void inc2() { synchronized(lock2) { c2++; } }
\end{verbatim}

\begin{tcolorbox}[tipbox,title=Reentrant synchronizace]
Vlákno může získat zámek, který \textbf{už vlastní} (např. synchronizovaná metoda volá jinou synchronizovanou metodu stejného objektu). Počítá se rekurzivně.
\end{tcolorbox}

\subsubsection{Guarded Blocks (wait/notify)}

\begin{verbatim}
public synchronized void guardedJoy() {
    while(!joy) {
        try {
            wait();  // uvolní zámek, čeká na notify
        } catch (InterruptedException e) {}
    }
    // pokračování...
}

public synchronized notifyJoy() {
    joy = true;
    notifyAll();  // probudí VŠECHNA čekající vlákna
}
\end{verbatim}

\begin{tcolorbox}[warnbox,title=Pravidla pro wait/notify]
\begin{itemize}
    \item Vždy volej \texttt{wait()} uvnitř \texttt{while} cyklu, ne \texttt{if} (spurious wakeup!)
    \item Vždy používej \texttt{synchronized} blok/metodu při volání \texttt{wait/notify}
    \item Preferuj \texttt{notifyAll()} před \texttt{notify()} (spravedlnost)
\end{itemize}
\end{tcolorbox}

\subsection{Problémy souběhu -- Deadlock, Starvation, Livelock}

\begin{center}
\begin{tabular}{>{\raggedright}p{2.5cm} >{\raggedright\arraybackslash}p{11cm}}
\toprule
\textbf{Problém} & \textbf{Popis} \\
\midrule
\textbf{Deadlock} & Dvě či více vláken navždy blokována, vzájemně čekají na zámky druhých. Příklad: vlákno A drží zámek 1, čeká na 2; vlákno B drží 2, čeká na 1. \\
\addlinespace
\textbf{Starvation} & Vlákno nemůže získat přístup ke zdrojům, protože jiná vlákna je neustále získávají. \\
\addlinespace
\textbf{Livelock} & Vlákna nejsou blokována, ale vzájemně reagují na akce druhých bez pokroku. Příklad: dva lidé se vyhýbají v chodbě, stále se posunují do stejné strany. \\
\bottomrule
\end{tabular}
\end{center}

\subsection{java.util.concurrent}

\subsubsection{Executor framework}

\begin{tcolorbox}[defbox,title=ExecutorService]
Odděluje \textbf{vytváření a management vláken} od aplikační logiky.
\begin{itemize}
    \item \texttt{Executor} -- základní rozhraní, metoda \texttt{execute(Runnable)}
    \item \texttt{ExecutorService} -- rozšiřuje o \texttt{submit()}, \texttt{shutdown()}, \texttt{awaitTermination()}
    \item \texttt{ScheduledExecutorService} -- plánované a opakované úlohy
\end{itemize}
\end{tcolorbox}

\begin{verbatim}
// Thread pool s fixním počtem vláken
ExecutorService executor = Executors.newFixedThreadPool(4);

// Odeslání úlohy
Future<Long> result = executor.submit(new Callable<Long>() {
    public Long call() { return computeValue(); }
});

// Získání výsledku (blokující)
Long value = result.get();
executor.shutdown();
\end{verbatim}

\begin{tcolorbox}[warnbox,title=Otázka ze vzorového testu!]
\textbf{ExecutorService}:
\begin{itemize}
    \item Určuje, \textbf{jakým způsobem a v jakém vlákně} budou úlohy provedeny
    \item Poskytuje metody produkující \texttt{Future} pro asynchronní návratové hodnoty
    \item Zajišťuje provedení předložených úloh
\end{itemize}
\textbf{Není pravda}, že neurčuje způsob provedení -- to je jeho hlavní účel!
\end{tcolorbox}

\subsubsection{Fork/Join framework}

\begin{itemize}
    \item Implementace \texttt{ExecutorService} pro \textbf{rekurzivní dělení} úloh
    \item \textbf{Work stealing} -- volné vlákno ``ukradne'' práci z fronty jiného
    \item Základ: \texttt{ForkJoinPool}, \texttt{RecursiveTask<T>} (vrací hodnotu), \texttt{RecursiveAction}
    \item Metoda \texttt{compute()}: pokud je úloha malá, proveď přímo; jinak rozděl a zavolej \texttt{invokeAll()}
\end{itemize}

\subsubsection{Konkurentní kolekce}

\begin{center}
\begin{tabular}{ll}
\toprule
\textbf{Třída/Interface} & \textbf{Účel} \\
\midrule
\texttt{BlockingQueue} & Fronta s blokujícím čtením/zápisem \\
\texttt{ArrayBlockingQueue} & Omezená kapacita \\
\texttt{ConcurrentHashMap} & Thread-safe HashMap \\
\texttt{ConcurrentSkipListMap} & Thread-safe TreeMap \\
\texttt{SynchronousQueue} & Předává prvky přímo (žádná kapacita) \\
\bottomrule
\end{tabular}
\end{center}

\subsection{Concurrency ve Swing}

\begin{tcolorbox}[defbox,title=Tři druhy vláken ve Swing]
\begin{enumerate}
    \item \textbf{Initial threads} -- počáteční vlákna (např. \texttt{main})
    \item \textbf{Event Dispatch Thread (EDT)} -- zpracovává události, manipuluje s GUI komponentami
    \item \textbf{Worker threads} -- časově náročné úlohy na pozadí
\end{enumerate}
\end{tcolorbox}

\begin{tcolorbox}[warnbox,title=Pravidlo EDT]
Swingové komponenty nejsou thread-safe! Veškerá manipulace s GUI musí probíhat v \textbf{EDT}. Používej:
\begin{verbatim}
SwingUtilities.invokeLater(() -> { /* GUI kód */ });
SwingUtilities.invokeAndWait(() -> { /* GUI kód */ }); // blokující
\end{verbatim}
\end{tcolorbox}

\subsubsection{SwingWorker}

\begin{verbatim}
SwingWorker<ImageIcon[], Void> worker = new SwingWorker<>() {
    @Override
    protected ImageIcon[] doInBackground() throws Exception {
        // provádí se v pracovním vlákně
        return loadImages();
    }
    
    @Override
    protected void done() {
        // provádí se v EDT po skončení
        try {
            ImageIcon[] result = get();
            // aktualizace GUI...
        } catch (Exception e) { ... }
    }
};
worker.execute();
\end{verbatim}

\begin{tcolorbox}[tipbox,title=Generické parametry SwingWorker]
\texttt{SwingWorker<T, V>}:
\begin{itemize}
    \item \texttt{T} -- návratový typ \texttt{doInBackground()} a \texttt{get()}
    \item \texttt{V} -- typ mezivýsledků pro \texttt{publish()}/\texttt{process()}
\end{itemize}
\end{tcolorbox}

%=============================================================================
\section{Shrnutí klíčových otázek ze vzorového testu}
%=============================================================================

\subsection{Odpovědi na viditelné otázky}

\begin{enumerate}[label=\textbf{Otázka \arabic*:}]
    \item \textbf{Který typ výjimek pro neočekávané stavy?} $\rightarrow$ \textbf{Nekontrolované výjimky} (RuntimeException)
    
    \item \textbf{K čemu slouží @FXML?} $\rightarrow$ \textbf{Propojení identifikace komponent} mezi kontrolerem a GUI
    
    \item \textbf{Nejvyšší v hierarchii JavaFX?} $\rightarrow$ \textbf{Stage}
    
    \item \textbf{Lze vyhodit výjimku v konstruktoru?} $\rightarrow$ \textbf{Ano} (bez omezení, není třeba uvádět v deklaraci pro unchecked; pro checked ano)
    
    \item \textbf{Co je defenzivní programování?} $\rightarrow$ \textbf{Předvídání chyb a mimořádných stavů}
    
    \item \textbf{Prostředky synchronizace:} $\rightarrow$ \textbf{Synchronizované metody, synchronizované bloky (zámky jsou implementační detail, konstruktory se nesynchronizují)}
    
    \item \textbf{Charakteristika ExecutorService:} $\rightarrow$ Určuje \textbf{způsob a vlákno} provedení, produkuje \texttt{Future}, zajišťuje provedení
    
    \item \textbf{Přístup vnitřní třídy k private členům?} $\rightarrow$ \textbf{Ano} -- vnitřní třídy mají přístup ke všem členům vnější, i private
    
    \item \textbf{Co je anonymní vnitřní třída?} $\rightarrow$ \textbf{Vnitřní třída, v jejíž deklaraci chybí hlavička} (jméno)
    
    \item \textbf{Co provádí Event Dispatch Thread?} $\rightarrow$ \textbf{Zpracovává systémovou frontu událostí týkajících se prvků GUI}
    
    \item \textbf{Vnitřní třída a dědičnost:} $\rightarrow$ \textbf{Vnitřní třída může být podtřídou vnější třídy} (není to pravidlo, ale je to možné)
    
    \item \textbf{Stavy vlákna:} $\rightarrow$ \textbf{New, Runnable, Blocked, Waiting, Timed waiting, Terminated}
    
    \item \textbf{Výhody genericity:} $\rightarrow$ Čitelnější kód, \textbf{kontroly při kompilaci} (ne přetypování!)
    
    \item \textbf{Omezení typových parametrů:} $\rightarrow$ \textbf{Omezuje množinu použitelných typových argumentů}
\end{enumerate}

%=============================================================================
\section{Rychlá reference -- často kladené otázky}
%=============================================================================

\begin{tcolorbox}[defbox,title=Rozdíly Swing vs. JavaFX]
\begin{tabular}{>{\raggedright}p{3.5cm} >{\raggedright}p{5cm} >{\raggedright\arraybackslash}p{5cm}}
\toprule
\textbf{Aspekt} & \textbf{Swing} & \textbf{JavaFX} \\
\midrule
Vlákno GUI & EDT & JavaFX Application Thread \\
Deklarativní GUI & Ne & FXML \\
CSS stylování & Omezené & Plná podpora \\
Scene Builder & Ne & Ano \\
3D grafika & Ne & Ano \\
WebView & Ne & Ano (WebKit) \\
Pracovní vlákna & SwingWorker & Task, Service \\
\bottomrule
\end{tabular}
\end{tcolorbox}

\begin{tcolorbox}[tipbox,title=Kontrolní seznam před testem]
\begin{enumerate}
    \item[$\square$] Rozdíl checked vs. unchecked exceptions
    \item[$\square$] Syntaxe try-catch-finally, vícenásobné catch
    \item[$\square$] throws klauzule, definice vlastních výjimek
    \item[$\square$] Účel assert (neplést s výjimkami!)
    \item[$\square$] Generické typy, wildcards, PECS pravidlo
    \item[$\square$] Type erasure -- co se děje při překladu
    \item[$\square$] Hierarchie JavaFX: Stage $>$ Scene $>$ Pane $>$ Node
    \item[$\square$] Účel @FXML anotace
    \item[$\square$] Vnořené třídy: statické vs. vnitřní, anonymní
    \item[$\square$] Stavy vlákna, životní cyklus
    \item[$\square$] Synchronizace: synchronized, ReentrantLock, atomické operace
    \item[$\square$] wait/notifyAll, Guarded Blocks
    \item[$\square$] Deadlock vs. Starvation vs. Livelock
    \item[$\square$] ExecutorService, ThreadPool, Future
    \item[$\square$] SwingWorker / Task + Service v JavaFX
    \item[$\square$] EDT vs. pracovní vlákna (proč oddělit GUI a výpočty)
\end{enumerate}
\end{tcolorbox}

\end{document}
